% !PRINTED 326
%<<CONT
des mittleren das Product der beiden äusseren abzieht, nämlich,
mit Benutzung der Relation $xf'(x) + yf'(y) = nf$:
\[
\frac{(Xy - Yx)^2}{nf^2}\left\{(n-1)H(Xy - Yx)^2 + [Xf'(x) + Yf'(y)]^2\right\},
\]
und man erhält also die Grenzwerthe von $X:Y$ aus der Gleichung

\weq{11}{[Xf'(x) + Yf'(y)]^2 + (n-1)(Xy - Yx)^2 H = 0,}

die in Bezug auf $X:Y$ quadratisch ist.

Da $H$ negativ ist, hat diese Gleichung zwei reelle Wurzeln,
die man aus

\weq{12}{Xf'(x) + Yf'(y) = \pm(Xy - Yx)\sqrt{(1-n)H}}

findet.

Diese Formel umfasst mehrere besondere Resultate, die man
erhält, wenn man die homogene Form verlässt.

Nehmen wir zunächst $y = 0$, also $x:y$ unendlich und
setzen $Y = 1$, so ergeben sich zwei Werthe $X_1$, $X_2$, zwischen
denen die Gesammtheit aller Wurzeln von $f(x)$ enthalten sind.
Setzen wir
\[
f(x, y) = a_0 x^n + a_1 x^{n-1} y + a_2 x^{n-2} y^2 + \cdots,
\]
so wird

\weq{13}{H = [2na_0a_2 - (n-1)a_1^2]x^{2n-4},}

und man erhält diese beiden Werthe aus

\weq{14}{na_0 X = -a_1 \pm \sqrt{(n-1)^2a_1^2 - 2n(n-1)a_0a_2}.}

Lässt man $x:y$ unbestimmt, kehrt aber zur inhomogenen
Form zurück, indem man $y = 1$, $Y = 1$ setzt und die Formeln
(7) und (9) des \S.\,106 anwendet, so erhält man aus (12)

\weq{15}{x - X = \frac{nf(x)}{f'(x) \pm \sqrt{(n-1)^2f'(x)^2 - n(n-1)f(x)f''(x)}},}

und diese Werthe haben folgende Bedeutung:

Liegt $x$ zwischen zwei Wurzeln $\alpha_1, \alpha_2$ von $f(x)$, so liegt von
den beiden durch (15) dargestellten Werthen von $X$ der eine
näher an $\alpha_1$, der andere näher an $\alpha_2$, aber beide noch innerhalb
des Intervalles $(\alpha_1, \alpha_2)$.

Ist aber $x$ grösser als die grösste, oder kleiner als die
kleinste Wurzel von $f(x)$, so liegen die sämmtlichen Wurzeln
zwischen den beiden Werthen von $X$, während $x$ ausserhalb liegt.
